\subsection{Relationship between energy terms and dynamical fields}\label{sec:correlation}

To assess how LEC terms covary with the dynamical fields shown in Sect.~\ref{sec:composites}, we computed Pearson linear correlations ($r$) between each canonical LEC term and each scalar feature derived from each field, and used Spearman rank correlation ($\rho$) as a monotonic-relationship check. Figure~\ref{fig:9_pearson_epall_by_field_type} displays only pairs with $|r| \geq 0.20$. In this subset ($126$ pairs), Pearson and Spearman signs agree in all cases (100\%), and only two pairs show $|\,|r|-|\rho|\,|>0.10$. Given $n=2729$--$2733$, all displayed correlations are highly significant statistically, but interpretation below is based primarily on effect size and physical consistency rather than on $p$-values alone.

A caveat is that these are contemporaneous, storm-centered, spatially aggregated features; therefore they diagnose covariability, not direction of causality. In addition, dependence among features (e.g., shared information between border, sector, and contrast metrics) reduces the number of effectively independent tests, so nominal significance levels should be interpreted conservatively in a field-significance sense \citep{livezey1983statistical,bretherton1999effective,wilks2016stippling}.

The $\mathrm{PV}_{850}$ panel contributes only $7/91$ pairs above threshold ($7.7\%$; maximum $|r|=0.27$), substantially fewer than all other analyzed fields, and is therefore not discussed in detail.

For $\mathrm{AdvT}_{850}$ (Fig.~\ref{fig:9_pearson_epall_by_field_type}a), the strongest links involve APE-related terms. For $C_a$, the largest magnitudes are in sector W $(r=-0.46,\,\rho=-0.55)$, sector N $(r=-0.45,\,\rho=-0.56)$, and domain mean $(r=-0.44,\,\rho=-0.50)$. For $A_e$, the largest values are in sector N $(r=-0.56,\,\rho=-0.64)$ and border N $(r=-0.54,\,\rho=-0.56)$. These signs are consistent with a stronger cold-advection signal on the cold side of the cyclone being associated with enhanced baroclinic conversion and larger eddy APE, in line with quasigeostrophic/baroclinic-growth arguments \citep{lorenz1955available,davies2015qgomega}.

For $BA_e$, $\mathrm{AdvT}_{850}$ correlations show a coherent east/west dipole: border E $(r=+0.28,\,\rho=+0.29)$, sector E $(r=+0.28,\,\rho=+0.30)$, and E--W contrast $(r=+0.33,\,\rho=+0.31)$ are positive, whereas border W $(r=-0.24,\,\rho=-0.20)$ and sector W $(r=-0.37,\,\rho=-0.28)$ are negative. This is consistent with stronger warm-side import and cold-side export of eddy APE across domain boundaries.

For $\mathrm{AFC}_{250}$ (Fig.~\ref{fig:9_pearson_epall_by_field_type}b), the dominant signal is the amplitude metric $|\mathrm{AFC}|$ (domain absolute mean), most strongly correlated with $K_e$ $(r=+0.53,\,\rho=+0.57)$ and $A_e$ $(r=+0.51,\,\rho=+0.56)$. The signed domain-mean AFC is positively associated with $C_a$ $(r=+0.43,\,\rho=+0.38)$ and negatively with $C_k$ $(r=-0.22,\,\rho=-0.23)$. These results are broadly consistent with downstream-development energetics, but should be interpreted as population-level co-variability rather than event-scale mechanism proof \citep{orlanski1991life,orlanski1993ageostrophic}.

The AFC-$BA_e$ pattern is also structured: border/sector W are positive [border W: $(r=+0.38,\,\rho=+0.37)$; sector W: $(r=+0.20,\,\rho=+0.19)$], while border E and E--W contrast are negative [border E: $(r=-0.29,\,\rho=-0.27)$; contrast E--W: $(r=-0.41,\,\rho=-0.40)$]. This supports an interpretation in which west-enhanced AFC configurations covary with stronger net $A_e$ import.

For $\mathrm{KE\_adv}_{250}$ (Fig.~\ref{fig:9_pearson_epall_by_field_type}c), the strongest individual pair is $A_e\times$sector W $(r=-0.44,\,\rho=-0.45)$. The amplitude metric $|\mathrm{KE\_adv}_{250}|$ is positively correlated with $A_e$ $(r=+0.43,\,\rho=+0.48)$, $C_a$ $(r=+0.34,\,\rho=+0.40)$, and $K_e$ $(r=+0.33,\,\rho=+0.40)$. For $C_k$, positive associations occur for domain mean and sectors N/S/E [domain mean: $(r=+0.33,\,\rho=+0.37)$; sector N: $(r=+0.26,\,\rho=+0.26)$; sector S: $(r=+0.23,\,\rho=+0.30)$; sector E: $(r=+0.27,\,\rho=+0.29)$]. Under the sign convention adopted in this manuscript ($C_k<0$ indicates mean-flow to eddy transfer), these positive correlations imply less negative $C_k$ where KE-advection metrics are larger.

An important dependence caveat is that $|\mathrm{AdvT}_{850}|$ and $|\mathrm{KE\_adv}_{250}|$ (domain absolute means) are themselves strongly correlated $(r=+0.61,\,\rho=+0.64;\,n=2733)$. Therefore, part of the KE-advection associations may reflect shared storm-intensity/advection variability rather than an independent upper-level pathway.

For $\mathrm{PV}_{200}$ (Fig.~\ref{fig:9_pearson_epall_by_field_type}d), the largest coefficients in the full figure are tied to E--W contrast: $G_e$ $(r=+0.59,\,\rho=+0.57)$, $A_e$ $(r=+0.57,\,\rho=+0.59)$, and $C_a$ $(r=+0.45,\,\rho=+0.47)$. In Southern Hemisphere convention, where $\mathrm{PV}_{200}$ is predominantly negative, a positive E--W contrast corresponds to less-negative PV to the east and more-negative PV to the west. This is consistent with a trough-west/ridge-east configuration associated with stronger diabatic generation and baroclinic energetics \citep{hoskins1985use,davis1991potential}.

Consistently, border/sector W PV values are negatively correlated with $C_a$ and $A_e$ [border W: $(r_{C_a}=-0.30,\,\rho=-0.31)$ and $(r_{A_e}=-0.44,\,\rho=-0.47)$; sector W: $(r_{C_a}=-0.24,\,\rho=-0.25)$ and $(r_{A_e}=-0.37,\,\rho=-0.39)$], indicating larger $C_a$ and $A_e$ under more cyclonic (more negative) western PV. By contrast, the isolated relation between border-W PV and $G_e$ is weak $(r=-0.05,\,\rho=-0.01)$, so this quantity should not be interpreted as a robust standalone predictor of $G_e$.

Overall, $\mathrm{AFC}_{250}$ and $\mathrm{KE\_adv}_{250}$ provide the largest fractions of pairs above threshold ($39.6\%$ and $38.5\%$), followed by $\mathrm{AdvT}_{850}$ ($33.0\%$) and $\mathrm{PV}_{200}$ ($19.8\%$), while $\mathrm{PV}_{850}$ remains sparse ($7.7\%$). The results support physically plausible population-level links between frontal/upper-level structure and LEC terms, but they should be treated as statistical constraints on mechanism hypotheses, not as direct causal attribution.